% ----------------------------------
% - Local Volatility Model
%   - Calibration
%     - Relation with implied volatilities
%     - Interpolation/extrapolation methods
%   - Experimental results
%     - Pricer (PDE + Monte Carlo)
%     - Calibration
% ----------------------------------

\chapter{Calibration of Local Volatility}

The local volatility model must be calibrated to market option prices to be usable.
Given quoted call prices $C^*(K_i, T_j)$ for a finite set of strikes and maturities,
we seek a function $\sigma(S,t)$, or equivalently $a(y,\tau) = \tfrac{1}{2}\sigma^2$,
such that the Dupire forward equation (\ref{eq:dupire_forward}) reproduces the observed prices.

We follow \cite{egger2005tikhonov} and treat the calibration as a Tikhonov regularization problem.

\section{Ill-posedness of the inverse pricing problem}

The calibration of the local volatility model is an ill-posed inverse problem,
because on one hand market data is noisy and sparse in maturity.
On the other hand the map from $a$ to option prices has an smoothing effect that makes its inverse discontinuous,
This is important, it means different volatility functions can produce almost the same option prices, especially away from the money.

We can see the latter in the Dupire equation \eqref{eq:dupire_general},
the gamma is in the denominator,
so in regions where gamma is small, like far OTM/ITM,
the local volatility is not well determined by the prices.

Also the Local Volatility model assumes that option prices are $C^{1,2}$ in $(T,K)$,
we are looking for a function $\sigma(S,t)$ that reproduces the market prices,
but PDE theory suggests that that the coefficient $\sigma$ should live in a more general space, like $H^1$ or $L^2$.
This arises the problem, where the solution is smoother than the coefficient.


\section{Change of variable - log-moneyness}

Similar to what we did in Section~\ref{sec:bs_pricing} for the Black-Scholes PDE,
we change variables to log-moneyness
\begin{equation*}
    y = \log(K/S_0), \quad \text{equivalently} \quad K = S_0\,e^{y},
\end{equation*}
removes this dependence and puts the PDE in a translation-invariant form,
just as the substitution $x = \log(S/K)$ reduced the Black-Scholes PDE to the heat
equation in Section~\ref{sec:bs_pricing}.
Define the time to maturity $\tau = T$,
the diffusion coefficient
\begin{equation*}
    a(y,\tau) = \tfrac{1}{2}\sigma^2(K,T),
\end{equation*}
and the normalized call price
\begin{equation*}
    u(y,\tau) = C(S_0; K, T)/S_0,
\end{equation*}
so that $C(S_0; K, T) = S_0\,u(y,\tau)$.
\footnote{We keep the explicit zero-order term. \cite{egger2005tikhonov} instead set
$u = e^{\int_t^T q\,ds}\,C$, which absorbs the dividend factor and yields the
equivalent equation $-u_\tau + a(u_{yy}-u_y) + (q-r)u_y = 0$ without a $-q\,u$ term;
the two normalizations coincide when $q = 0$, as in Example~1 below. We implemented and
tested both on the market surface, where $q \neq 0$, and neither outperformed the other;
see Appendix~\ref{subsec:discounted_normalization}.}

We now carry the Dupire forward equation~\eqref{eq:dupire_forward} into the
$(y,\tau)$ coordinates. Since $y = \log(K/S_0)$ we have $\partial y/\partial K = 1/K$,
and because $\tau = T$ the maturity derivative is unchanged.
Applying the chain rule to $C = S_0\,u(y,\tau)$ gives the three derivatives that appear
in~\eqref{eq:dupire_forward}:
\begin{equation*}
\begin{aligned}
    \frac{\partial C}{\partial T}
        &= S_0\,u_\tau, \\
    \frac{\partial C}{\partial K}
        &= S_0\,u_y\,\frac{\partial y}{\partial K}
         = \frac{S_0}{K}\,u_y, \\
    \frac{\partial^2 C}{\partial K^2}
        &= \frac{\partial}{\partial K}\!\left(\frac{S_0}{K}\,u_y\right)
         = \frac{S_0}{K}\,u_{yy}\,\frac{\partial y}{\partial K} - \frac{S_0}{K^2}\,u_y
         = \frac{S_0}{K^2}\,(u_{yy} - u_y).
\end{aligned}
\end{equation*}
The convenient feature of log-moneyness is that the first-order operator $K\,\partial_K$
and the second-order operator $K^2\,\partial_{KK}$ both become constant-coefficient
operators in $y$:
\begin{equation*}
    K\,\frac{\partial C}{\partial K} = S_0\,u_y,
    \qquad
    K^2\,\frac{\partial^2 C}{\partial K^2} = S_0\,(u_{yy} - u_y),
\end{equation*}
which is exactly what makes the equation translation-invariant in $y$.
Substituting these into \eqref{eq:dupire_forward} and using $a = \tfrac12\sigma^2$,
\begin{equation*}
    S_0\,u_\tau + (r-q)\,S_0\,u_y - a\,S_0\,(u_{yy} - u_y) = -q\,S_0\,u.
\end{equation*}
Dividing through by $S_0$ and solving for $u_\tau$ yields the forward equation in
log-moneyness form,
\begin{equation}\label{eq:inverse_pde_1}
    u_\tau = a(y,\tau)\,(u_{yy} - u_y) - (r-q)\,u_y - q\,u,
\end{equation}
on the truncated domain $Q_{M,T} = (-M, M) \times (0, T]$.
The initial condition follows from the same normalization applied to the
call payoff $C(K,0) = (S_0 - K)^+$:
\begin{equation*}
    u(y,0) = \frac{(S_0 - K)^+}{S_0} = \left(1 - \frac{K}{S_0}\right)^{\!+} = (1 - e^{y})^+,
\end{equation*}
that is,
\begin{equation}\label{eq:inverse_pde_3}
    u(y, 0) = \max(1 - e^y, 0)
\end{equation}
and Dirichlet boundary conditions
\begin{equation}\label{eq:inverse_pde_2}
    u(-M, \tau) = \Psi_1(\tau), \qquad u(M, \tau) = \Psi_2(\tau).
\end{equation}
The boundary values are determined from the Black-Scholes formula evaluated
at the boundary strikes $K = S_0 e^{\pm M}$:
as $M \to \infty$, the left boundary ($K \to 0$) gives the intrinsic value
$\Psi_1(\tau) \to e^{-q\tau}$, and the right boundary ($K \to \infty$) gives $\Psi_2(\tau) \to 0$.


\section{Parameter-to-solution map}

For each admissible diffusion coefficient $a$,
let $u(a)$ denote the solution to \eqref{eq:inverse_pde_1} - \eqref{eq:inverse_pde_3}.
Fix an \emph{a priori} reference $a^* \in L^\infty(\mathbb{R})$ with $0 < \underline{a} \le a^* \le \bar{a}$,
where $\underline{a}, \bar{a}$ are self-imposed bounds on the volatility,
and define the set of admissible parameters
\begin{equation}\label{eq:admissible_set}
    \mathcal{K}(a^*) := \bigl\{ a \in a^* + H^1(Q) : \underline{a} \le a \le \bar{a} \bigr\}.
\end{equation}
Here $Q = \Omega \times (0,T)$ is the space-time domain and the span space $a^* + H^1(Q)$ is taken such as
$\nabla a^* \in [L^2(Q)]^2$ (Definition~2.1 of \cite{egger2005tikhonov}).

Two features of $H^1$ matter.
As we will present afterwards in \eqref{eq:variational_forward_eq},
the coefficient enters the equation through the flux term,
so it must have a weak derivative in $L^2$.
And also the $H^1$ penalty makes the sublevel sets of the Tikhonov functional bounded in $H^1$,
hence weakly compact, which is what yields a minimizer.
In the synthetic examples below the volatility is time independent, $a = a(y)$,
and the penalty reduces to the spatial norm
\begin{equation*}
    \int_{-M}^{M} (a^2 + a_y^2)\,dy,
\end{equation*}
on the bounded interval $H^1$ then also embeds into $C[-M,M]$,
so the bounds $\underline{a} \le a \le \bar{a}$ hold pointwise.

The \emph{parameter-to-solution map} $F = F_{a_0}$ is defined as
\begin{equation}\label{eq:param_solution_map}
    F(a) = u(a) - u(a_0),
\end{equation}
where $a_0 \in \mathbb{R}$ with $\underline{a} \le a_0 \le \bar{a}$ is a reference constant
and $u(a_0)$ can be computed from the Black-Scholes formula.
A key regularity result established in \cite{egger2005tikhonov} is the following.

\begin{theorem}[\cite{egger2005tikhonov}, Theorem 2.1]
There exists $p^* > 2$ such that $F\colon \mathcal{K}(a^*) \to W_p^{2,1}(Q)$
is continuous and compact for $2 \le p < p^*$, and weakly sequentially continuous.
\end{theorem}

The compactness and weak closeness of $F$ is the mathematical source of the ill-posedness.
As already mentioned it reflects the smoothing action of the pricing equation,
the solution of a parabolic problem is more regular than its coefficient,
and the solution space $W_p^{2,1}(Q)$ embeds compactly in the observation space.
A compact operator between infinite dimensional spaces cannot have a continuous inverse,
so arbitrarily small changes in the observed prices can correspond to order one changes in $a$.
In practice very different volatility profiles reproduce almost the same option prices, especially away from the money,
and any direct inversion amplifies the noise.

Controlling this instability is the purpose of the Tikhonov penalty.


\section{Tikhonov functional}

Lets assume that there exists $a^\dagger \in \mathcal{K}(a^*)$ such that $u^\dagger = u(a^\dagger)$ reproduces the 'true' prices $u^\dagger$ exactly.
And let $u^\delta$ denote the noisy observed option prices with a bound known a priori,
\begin{equation*}
    \|u^\delta - u^\dagger\|_{L_2(Q)} \le \delta.
\end{equation*}
Following \cite{egger2005tikhonov} we stabilize the inverse problem by minimizing the
\emph{Tikhonov functional}
\begin{equation}\label{eq:tikhonov_functional}
    J^\beta(a) = \|u(a) - u^\delta\|^2_{L_2} + \beta\,\|a - a^*\|^2_{H^1},
\end{equation}
where $\beta > 0$ is the regularization parameter,
and $a^*$ is the \emph{a priori} guess that encodes any structural knowledge about the volatility (e.g., smoothness or smile shape).
The first term penalizes any mismatch with the data,
the second term penalizes the deviation of $a$ from $a^*$ in $H^1$,
which promotes smoother reconstructions.

The penalty norm in $H^1$ is given by
\begin{equation*}
    \|a - a^*\|^2_{H^1} 
        = \int_{-M}^{M} \bigl[(a-a^*)^2 + (\partial_y(a-a^*))^2\bigr]\,dy.
\end{equation*}

\begin{proposition}[\cite{egger2005tikhonov}, Propositions 3.1 and 3.2]
\label{prop:tikhonov_convergence}
For any $\beta > 0$ a minimizer $a^\delta_\beta$ of $J^\beta$ exists.
If $\delta \to 0$ and $\beta(\delta) \to 0$ with $\delta^2/\beta(\delta) \to 0$,
then every sequence of minimizers $\{a^{\delta_k}_{\beta(\delta_k)}\}$
has a convergent subsequence, and every limit point is an $a^*$-minimum-norm solution.
\end{proposition}

% \begin{proof}[Idea of the proof]
% Both statements are the standard results for Tikhonov regularization of nonlinear ill-posed
% problems~\cite{engl_hanke_neubauer}, which Theorem~2.1 makes applicable here. For existence,
% a minimizing sequence $\{a_n\}$ is bounded in $H^1$ by the penalty term, so it has a weakly
% convergent subsequence $a_n \rightharpoonup \bar{a}$. The constraint set is convex and closed,
% hence weakly closed, so $\bar{a} \in \mathcal{K}(a^*)$. Weak sequential continuity of $F$
% (Theorem~2.1) and weak lower semicontinuity of the $H^1$ norm give
% $J^\beta(\bar{a}) \le \liminf_n J^\beta(a_n)$, so $\bar{a}$ is a minimizer. For the vanishing
% noise limit, the choice $\delta^2/\beta(\delta) \to 0$ keeps the data and penalty terms
% balanced, so the minimizers stay bounded and every weak limit point solves the noise free
% problem with least $H^1$ distance to $a^*$. We refer to~\cite{engl_hanke_neubauer} for the
% general theory and to~\cite{egger2005tikhonov} for its application here.
% \end{proof}

An important remark is done in \cite{egger2005tikhonov}.
Usually in the analysis of the inverse problems,
in particular with the Tikhonov regularization approach,
is assumed that the observations of the solutions of the PDE are continuous,
in our case we would like to have as data the prices of the options in any strike or maturity,
as a continuous function.
But in practice we only have a finite number of prices, so the data is discrete.
This makes the convergence theory more complicated,
in particular the convergence to some $a^*-minimum-norm$ solution is guaranteed,
but its uniqueness is not.

\section{Gradient computation via the adjoint method}

Before tackling the convergence of the minimization method,
we need to compute the gradient of the Tikhonov functional $J^\beta$,
in particular we need its Gâteaux derivative.
The \emph{directional derivative} of the data-fidelity term at $a$ in direction $h$ is
\begin{equation*}
    \frac{d}{d\varepsilon}\|u(a+\varepsilon h) - u^\delta\|^2_{L_2}\Big|_{\varepsilon=0}
    = 2\bigl(u(a) - u^\delta,\; u'(a)[h]\bigr)_{L_2},
\end{equation*}
where $u'(a)[h] = w(\cdot, T_0)$ is the solution of the \emph{linearized} forward equation,
obtained by differentiating~\eqref{eq:inverse_pde_1} in the direction $h$,
\begin{equation}\label{eq:gateaux_derivative}
    -w_\tau + a(w_{yy} - w_y) - (r-q)\,w_y - q\,w = -h\,(u_{yy} - u_y),
    \quad w(y,0)=0, \quad w\big|_{\partial\Omega}=0.
\end{equation}
Rather than solving \eqref{eq:gateaux_derivative} for every direction $h$,
the gradient is evaluated efficiently via a single \emph{adjoint solve}.
Let $V(y,\tau)$ solve the backward parabolic problem
\begin{equation}\label{eq:adjoint_pde}
    V_\tau + (aV)_{yy} + (aV)_y + (r-q)V_y - qV = 0,
    \qquad V(y, T_0) = u(a)(y, T_0) - u^\delta(y),
\end{equation}
with homogeneous Dirichlet conditions at $y = \pm M$ and integrated backward in $\tau$.
Then the Gâteaux derivative satisfies \cite{egger2005tikhonov}
\begin{equation*}
    \bigl(u'(a)[h],\, u(a)-u^\delta\bigr)_{L_2}
    = \int_0^{T_0}\!\!\int_{-M}^{M} h(y)\,(u_{yy}-u_y)\,V\,dy\,d\tau.
\end{equation*}
The regularization term, is an explicit functional of $a$,
so its Gâteaux derivative needs no adjoint,
we find it by differentiating the quadratic form directly gives
\begin{equation*}
    \frac{d}{d\varepsilon}\|a+\varepsilon h - a^*\|^2_{H^1}\Big|_{\varepsilon=0}
    = 2(a-a^*,h)_{H^1}
    = 2\int_{-M}^{M}\bigl[(a-a^*)h + \partial_y(a-a^*)\partial_y h\bigr]dy.
\end{equation*}
Adding this to the data term above, the total gradient direction (in $L^2$ sense) is
\begin{equation}\label{eq:tikhonov_gradient}
    \langle \nabla J^\beta(a),\, h \rangle
    = 2\int_0^{T_0}\!\!\int_{-M}^{M} h\,(u_{yy}-u_y)\,V\,dy\,d\tau
      + 2\beta\int_{-M}^{M}\bigl[(a-a^*)h + \partial_y(a-a^*)\partial_y h\bigr]dy.
\end{equation}
\section{Convergence rates and the source condition}\label{sec:convergence_rates}

The adjoint calculation of the previous section computes the action of $F'(a)^*$ on the residual.
Writing the data discrepancy as $\|F(a) - u^\delta\|^2$ and using
\begin{equation*}
    \langle \nabla J_d(a),\, h \rangle
    = 2\bigl(F'(a)h,\; F(a) - u^\delta\bigr)
    = 2\bigl(h,\; F'(a)^*(F(a) - u^\delta)\bigr),
\end{equation*}
the solution $V$ of \eqref{eq:adjoint_pde} represents $F'(a)^*(F(a) - u^\delta)$ (Remark 4.2 of \cite{egger2005tikhonov}).
The same adjoint controls the speed of convergence.

The convergence results exposed in proposition \ref{prop:tikhonov_convergence} can be quite slow,
they give no rate of convergence.
Without a further assumptions we can no derive such a rate of convergence.
The general theory \cite{engl_hanke_neubauer} provides a rate under a source condition on the true coefficient,
\begin{equation}\label{eq:source_condition}
    a^* - a^\dagger = F'(a^\dagger)^*\,w,
\end{equation}
holds for small $\|w\|$.
And $F$ is Frechet differentiable and $F'$ is Lipschitz continuous in a neighborhood of $a^\dagger$.
Though in \cite{egger2005tikhonov} is mentioned that $F$ just needs to be Gateaux differentiable and Lipschitz continuous in some directions,
so we will be able to use the Gateaux derivative we just computed.

Proposition 4.1 of \cite{egger2005tikhonov} and $\|w\|$ small enough,
the choice $\beta \sim \delta$ gives convergence rates
\begin{equation}\label{eq:convergence_rate}
    \|a^\delta_\beta - a^\dagger\| = O(\sqrt{\delta}),
    \qquad
    \|u(a^\delta_\beta) - u^\delta\| = O(\delta).
\end{equation}

The abstract condition \eqref{eq:source_condition} is hard to check directly,
so Egger and Engl derive interpretable sufficient conditions by restricting to a single maturity and a time independent volatility,
$a(\tau, y) = a(y)$.
In that setting they prove two rate results of the form \eqref{eq:convergence_rate}.
First, Theorem 4.1 treats observations of prices,
with the penalty strengthened to the $H^2$ norm.
Second, Theorem 4.2 keeps the $H^1$ penalty but assumes that the derivative of the option price is also observable,
which is realistic because $C_K = e^{-y}u_y$ is the price of a digital call and is quoted on the market.

In both, the source function must be smooth and decay like a Gaussian at both ends,
with $g = \partial_y^4 w - \partial_y^2 w + w$ for the first result and $g = -w_{yy} + w$ for the second,
the requirement is $|g|, |g_y|, |g_{yy}| = O(e^{-B' y^2})$ as $|y| \to \infty$.

This decay condition is the precise form of the wing dynamics,
and Egger and Engl argue the same in their concluding remarks,
option prices in a sense reflect the market's expectation of future volatility,
which is reasonably smooth,
but for deep in-the-money/ out-of-the-money options,
volatility is too noisy given the data,
so that structural behavior must come from the prior $a^*$ \cite{egger2005tikhonov}.
The condition requires $a^\dagger - a^*$ to be small and smooth in the wings,
which is why the priors in the experiments below are chosen to flatten away from the money.

Two points carry over to the rest of the chapter, and neither is proved by the theorems above.
The functional we minimize keeps the $H^1$ penalty with price data only,
which is the setting of the stability and convergence results but not literally of either rate theorem,
since those need the stronger $H^2$ penalty or the added derivative data.
And the rate analysis is proved only for a time independent volatility,
so the genuinely time dependent surface $a(y,\tau)$ we calibrate in Section~\ref{sec:market_surface} lies outside its scope.

\section{Choice of $\beta$ and stopping criterion}\label{sec:beta_choice}

In the numerical experiments of \cite{egger2005tikhonov} the regularization parameter
is chosen by a decreasing sequence
\begin{equation*}
    \beta_0 = 10^{2}\delta > \beta_1 > \cdots > \beta_n = 10^{-3}\delta,
\end{equation*}
minimizing $J^{\beta_i}$ sequentially,
starting each minimization step from the previous minimizer solution.

The sequence is stopped when the Morozov discrepancy principle
\begin{equation*}
    \|u(a^\delta_\beta) - u^\delta\| \le \tau\delta
\end{equation*}
is satisfied for some $\tau > 1$.
In \cite{egger2005tikhonov} they mention that after testing across the noise levels $\delta \in [10^{-5}, 10^{-1}]$,
the choice $\tau = 1.5$ could be met at typical values $\beta \sim 10^{-2}\delta$.
For exact data (i.e. no noise, $\delta = 0$) a single fixed $\beta$ is used.

\section{Variational forms and discretization}

The forward PDE \eqref{eq:inverse_pde_1} is discretized in space with
$P_1$ finite elements on a mesh of $[-M, M]$ and in time with the
implicit Euler scheme.
We can find the variational form by multiplying by a test function $v \in H^1_0(-M,M)$ and integrating by parts gives
\begin{equation}\label{eq:variational_forward_eq}
    \frac{d}{d\tau}\int_{-M}^{M} u\,v\,dy
    = -\int_{-M}^{M} a\,u_y\,v_y\,dy
      -\int_{-M}^{M}\partial_y a\,u_y\,v\,dy
      -\int_{-M}^{M}(a + r - q)\,u_y\,v\,dy
      -\int_{-M}^{M} q\,u\,v\,dy.
\end{equation}
The adjoint equation \eqref{eq:adjoint_pde} is run backward in $\tau$,
again multiplying by a test function $v \in H^1_0(-M,M)$ and integrating by parts gives its weak form (with $m = V$)
\begin{equation}\label{eq:variational_adjoint_eq}
    \frac{d}{d\tau}\int_{-M}^{M} m\,v\,dy
    = \int_{-M}^{M}(am)_y\,v_y\,dy
      +\int_{-M}^{M} a\,m\,v_y\,dy
      +\int_{-M}^{M}(r-q)\,m\,v_y\,dy
      +\int_{-M}^{M} q\,m\,v\,dy.
\end{equation}

\subsection{Time discretization}

We discretize of time by $N$ arbitrary small steps $\tau_0, \tau_1, \ldots, \tau_N$ with the implicit Euler scheme,
which is defined by
\begin{equation}\label{eq:backward_euler}
    \frac{1}{\Delta\tau}\int_{-M}^{M}\bigl(u^n - u^{n-1}\bigr)v\,dy
    = \,\phi\bigl(a^n;\,u^n,\,v\bigr)
    \qquad\text{for all } v \in V_h \cap H^1_0(-M,M),
\end{equation}
where $u^n = u(\cdot,\tau_n)$ and $\phi(a^n;u^n,v)$ is the right-hand side of the variational form \eqref{eq:variational_forward_eq}.
This gives a linear system for each step,
is unconditionally stable for this class of parabolic operators, see \cite{thomee_galerkin}.

\subsection{A time-dependent coefficient}\label{sec:time_dependent_a}

The admissible set \eqref{eq:admissible_set} is stated on the space $Q = \Omega\times(0,T)$,
so the theory shown in Proposition \ref{prop:tikhonov_convergence} covers time-dependent coefficients.
However, as previously stated,
in practice option prices are observed at finitely many maturities $\tau_1 < \cdots < \tau_N$,
and anything finer than the gaps between quoted maturities is not identifiable.
We therefore choose to model $a$ as piecewise constant in time,
\begin{equation}\label{eq:coefficient_slices}
    a(y,\tau) = \sum_{n=1}^{K} a_n(y)\,\chi_{(\tau_{n-1},\,\tau_n]}(\tau),
    \qquad a_n \in H^1(-M,M), \quad \underline{a} \le a_n \le \bar a,
\end{equation}
with $\tau_0 = 0$.
The unknown is the stack of spatial slices $(a_1, \dots, a_N)$,
in the case $N = 1$ recovers the time-independent setting of \cite{egger2005tikhonov}.
The Tikhonov functional is adapted accordingly
\begin{equation}\label{eq:surface_functional}
    J(a) = \sum_{j} \bigl\| u(a)(\cdot,\tau_j) - u^\delta_j \bigr\|^2_{L_2(\Omega_j)}
         + \beta_y \sum_{n} \| a_n - a^*_n \|^2_{H^1}
         + \beta_\tau \sum_{n} \| a_n - a_{n-1} \|^2_{L_2},
\end{equation}
where $\Omega_j$ is the observed strike band at maturity $\tau_j$.
The $H^1$ term smooths in strike,
the new $L_2$ term penalizes jumps between consecutive slices.

\section{Minimization Algorithm}

We minimize the discretized Tikhonov functional with L-BFGS-B \cite{byrd_lbfgsb},
a quasi-Newton method that handles the bound constraints $0 < \underline{a} \le a_k \le \bar a$.
Each iteration requires one forward solve of \eqref{eq:backward_euler} and one backward sweep of the adjoint equation \eqref{eq:variational_adjoint_eq},
which together assemble the gradient \eqref{eq:tikhonov_gradient} at the cost of two PDE solves per iteration.
Independent of the number of degrees of freedom.

\subsection{Verification of the discrete adjoint gradient}\label{sec:adjoint_verification}

The adjoint gradient is verified by two tests.

\paragraph{Tangent-linear/adjoint duality.}
Write the data-fidelity term as
\begin{equation*}
    J_d(a)=\|u(a)-u^\delta\|^2_{L_2}
\end{equation*}
with residual $r = u(a)-u^\delta$.
The gradient returned by the adjoint solve should satisfy
\begin{equation*}
    \bigl\langle \nabla J_d(a),\, h\bigr\rangle_{\mathbb{R}^n}
    = 2\bigl(r,\; u'(a)[h]\bigr)_{L_2}
\end{equation*}
for every direction $h$,
with the right-hand side computed separately from the linearized solve \eqref{eq:gateaux_derivative}.
This equality is exact, with no finite-difference truncation,
so if it holds the discrete adjoint is the exact transpose of the linearized forward operator.
For a random $h$ on the Example~1 grid the two sides differ by $2.3\times10^{-15}$,
i.e. at the level of rounding error.

\paragraph{Taylor remainder test.}
When the gradient is correct, the first-order Taylor remainder of the full functional decays
one order faster in $t$ than the zeroth-order one
\begin{equation*}
    R_0(t) = \bigl|J(a+th) - J(a)\bigr| = O(t),
    \qquad
    R_1(t) = \bigl|J(a+th) - J(a) - t\,\langle\nabla J(a), h\rangle\bigr| = O(t^2).
\end{equation*}
Table \ref{tab:taylor_test} lists $R_0$ and $R_1$ as the step $t$ is reduced geometrically.
The remainder $R_1$ falls at rate $\approx 2$, as it should when $\nabla J$ is exact.

\begin{table}[h]
    \centering
    \begin{tabular}{c c c c}
        \toprule
        $t$ & $R_0(t)$ & $R_1(t)$ & rate of $R_1$ \\
        \midrule
        $10^{-2}$ & $1.32\times10^{-6}$ & $5.13\times10^{-6}$ & - \\
        $10^{-3}$ & $3.28\times10^{-7}$ & $5.27\times10^{-8}$ & $1.99$ \\
        $10^{-4}$ & $3.76\times10^{-8}$ & $5.29\times10^{-10}$ & $2.00$ \\
        $10^{-5}$ & $3.80\times10^{-9}$ & $5.29\times10^{-12}$ & $2.00$ \\
        \bottomrule
    \end{tabular}
    \caption{
        Taylor remainder test for the discrete adjoint gradient (Example 1 setup \ref{sec:egger_ex1}).
        The zeroth-order remainder $R_0$ decays at first order in $t$ and the
        first-order remainder $R_1$ at second order (rate $\approx 2$),
        which is what an exact gradient produces.
        The rate column is $\log_{10}\!\bigl(R_1(10\,t)/R_1(t)\bigr)$.
    }
    \label{tab:taylor_test}
\end{table}

\section{Numerical experiment - Example 1}\label{sec:egger_ex1}

We reproduce \emph{Example 1} of \cite{egger2005tikhonov},
where we try to recover a constant volatility on $y \in [-3,3]$ ($M = 3$), $\tau \in [0,1]$,
with $r = q = 0$ and spot $S_0 = 1000$.
\footnote{
    The paper's text states $S_0 = 100$,
    but its tabulated option values (e.g.\ a call worth $439$ at $K = 600$)
    and test~(D) (``an underlying worth $1000\$$'')
    are only consistent with $S_0 = 1000$, we use the latter.
    The recovered $a(y)$ is scale-invariant,
    but the Tikhonov balance is not,
    the data term scales as $S_0^2$ while the penalty does not,
    so reproducing the paper's $\beta$ requires its $S_0$.
}
The true parameter is
\[
    a^\dagger = 0.15 \quad \Longleftrightarrow \quad \sigma^\dagger = \sqrt{2 \times 0.15} \approx 0.548.
\]

For part \emph{(A)} of the example,
the data $u^\delta$ are the exact option prices at $\tau = T = 1$,
generated from the Black-Scholes formula.
The initial guess and $a^*$ are both set to
\[
    a^*_1(y) = 0.15 - 0.05\,\operatorname{erf}(-y^2),
\]
which agrees with the true value at the money and the wings flatten out ($\partial_y a^* \to 0$,
so $a^* \to 0.20$) for large $|y|$, the asymptotic behavior invoked in
\cite{egger2005tikhonov} for the convergence-rate theorems 4.1-4.2.
The regularization parameter is fixed at $\beta = 10^{-6}$.

Subsequent parts increase the difficulty of the problem,
in part \emph{B} consists on taking 20 strikes only.
For part \emph{C} the prior is set to the deliberately poor constant $a^* = 0.15$.
Part \emph{D} adds $0.1\%$ noise to the continuous data.
And finally part \emph{BD} is both incomplete and noisy.

Our results are shown in following table \ref{tab:example_1}
\begin{table}[h]
    \centering
    \begin{tabular}{c c c c c c c}
        \toprule
        Strike & True Value & A & B & C & D & BD \\
        \midrule
        600 & 0.1500 & 0.1493 & 0.1500 & 0.1519 & 0.1378 & 0.1245 \\
        700 & 0.1500 & 0.1498 & 0.1503 & 0.1491 & 0.1386 & 0.1753 \\
        800 & 0.1500 & 0.1504 & 0.1499 & 0.1503 & 0.1364 & 0.1417 \\
        900 & 0.1500 & 0.1494 & 0.1499 & 0.1514 & 0.1372 & 0.1589 \\
        1000 & 0.1500 & 0.1499 & 0.1499 & 0.1515 & 0.1490 & 0.1455 \\
        1100 & 0.1500 & 0.1498 & 0.1500 & 0.1511 & 0.1581 & 0.1464 \\
        1200 & 0.1500 & 0.1496 & 0.1500 & 0.1502 & 0.1461 & 0.1595 \\
        1300 & 0.1500 & 0.1507 & 0.1502 & 0.1491 & 0.1506 & 0.1397 \\
        1400 & 0.1500 & 0.1497 & 0.1498 & 0.1504 & 0.1489 & 0.1522 \\
        1500 & 0.1500 & 0.1497 & 0.1502 & 0.1507 & 0.1465 & 0.1587 \\
        1600 & 0.1500 & 0.1508 & 0.1501 & 0.1491 & 0.1543 & 0.1500 \\
        1700 & 0.1500 & 0.1501 & 0.1496 & 0.1499 & 0.1496 & 0.1449 \\
        1800 & 0.1500 & 0.1487 & 0.1504 & 0.1511 & 0.1461 & 0.1463 \\
        \bottomrule
    \end{tabular}
    \caption{
        Example 1 of \cite{egger2005tikhonov},
        recovering $a(y)$ at the strikes $K = 600$ to $K = 1800$ for the five test cases.
    }
    \label{tab:example_1}
\end{table}

Figures \ref{fig:egger_ex1_B}, \ref{fig:egger_ex1_C} and \ref{fig:egger_ex1_BD},
illustrate how good of a fit we managed to get.
As shown the observed region is much more narrow than the domain,
we manage to recover a good approximation of the true value in the observed region but on the wings looks like the prior dominates the solution,
this is quite clear if we compare figures \ref{fig:egger_ex1_B} and \ref{fig:egger_ex1_C}.
\begin{figure}[h!]
    \centering
    \includegraphics[width=\textwidth]{egger_example1_caseB}
    \caption{
        Example 1 case \emph{B} of \cite{egger2005tikhonov},
        recovering the constant $a(y) = 0.15$ using only 20 strikes.
    }
    \label{fig:egger_ex1_B}
\end{figure}
\begin{figure}[h!]
    \centering
    \includegraphics[width=\textwidth]{egger_example1_caseC}
    \caption{
        Example 1 case \emph{C} of \cite{egger2005tikhonov},
        recovering the constant $a(y) = 0.15$ from a poor constant prior $a^* = 0.10$.
    }
    \label{fig:egger_ex1_C}
\end{figure}

While in \emph{D} and \emph{BD} we can see that the noise in the data makes the recovery of the true value more difficult,
we interpreted that the $0.1\%$ noise is applied to the call prices,
thus deep ITM options would have a bigger shock in absolute terms than deep OTM options.
We don't know if that was \emph{Egger's} and \emph{Engl's} approach,
as they mention that the noise would be $ 1\$ $, as the underlying is $1000\$ $,
but that would probably make some deep OTM options to have negative prices,
which is not possible,
thus recovering the true value would be impossible in those cases.

\begin{figure}[h!]
    \centering
    \includegraphics[width=\textwidth]{egger_example1_caseBD}
    \caption{
        Example 1 case \emph{BD} of \cite{egger2005tikhonov},
        recovering the constant $a(y) = 0.15$ from incomplete and noisy data.
    }
    \label{fig:egger_ex1_BD}
\end{figure}

\subsection{Numerical experiment - Example 2}

Example 2 of \cite{egger2005tikhonov} tries to show the importance of regularization.
The set up is similar to example 1,
but its true coefficient $a^\dagger(y)$ is the oscillatory profile near ATM
\begin{equation}\label{eq:egger_ex2_true}
    a^\dagger(y) = 0.15 + 0.05\,e^{-(y+0.3)^2}\sin(2\pi y) + 0.05\,\operatorname{erf}(20 y).
\end{equation}
To avoid an \emph{inverse crime} we generate the synthetic data on a wider domain $[-5,5]$,
with a four-times finer mesh and time step,
and interpolate them onto the coarser reconstruction grid,
so the data do not come from the same discretized operator that is later inverted.
We do it because \cite{egger2005tikhonov} mentions to use a different time-integration scheme for the same reason,
but does not mention the details,
in some of the cases looks like we are getting better reconstruction results and that might be the reason.

We run the paper's five cases.
All use the \emph{a priori} guess
\begin{equation*}
    a^*_1(y) = 0.15 + 0.05 \, \operatorname{erf}(2y),
\end{equation*}
which has the correct wings, except case C,
\begin{itemize}
    \item[(A)] complete, exact data;
    \item[(B)] incomplete data, only $20$ strikes observed;
    \item[(C)] complete, exact data with the poor constant prior $a^* = 0.15$;
    \item[(D)] complete data with $0.1\%$ noise;
    \item[(BD)] incomplete and noisy.
\end{itemize}
For the exact data cases we use a single small $\beta = 10^{-6}$.
Whereas for the noisy cases D and BD select $\beta$ with the decreasing $\beta$ schedule and Morozov discrepancy principle described above,
starting from $\beta_0 = 10^2\delta$ and warm-starting each minimization from the previous one.

Figure \ref{fig:egger_ex2} shows the recovered coefficient for the five cases.
Most exemplifying one is \emph{C},
where we can see the importance of the prior,
the poor constant prior $a^* = 0.15$ is not able to recover the oscillatory true value $a^\dagger(y)$,
specially in the wings where the data is weak and the prior dominates the solution.
\begin{figure}[h!]
    \centering
    \includegraphics[width=\textwidth]{egger_example2_recovery}
    \caption{Example~2 of \cite{egger2005tikhonov}: recovery of the oscillatory local
             variance $a^\dagger(y)$~\eqref{eq:egger_ex2_true} (red dashed) for the five test
             cases. Blue: recovered $a(y)$; grey dotted: prior $a^*$; black dots (B, BD):
             observed strikes; shaded band: the strikes $K = 600$--$1800$ tabulated in the
             paper.}
    \label{fig:egger_ex2}
\end{figure}

\section{Numerical experiment - a local volatility surface from market data}\label{sec:market_surface}

Examples 1 and 2 use synthetic data and a coefficient that is constant in time. We now
calibrate a full surface $a(y,\tau)$ from real market quotes. This goes past the setting of
Section~\ref{sec:convergence_rates}, where the rate theory assumes a time independent
coefficient: here the unknown is genuinely time dependent, so the surface fit is a numerical
extension rather than a covered case.

The data are SPY European option quotes from 30 April 2026, with spot $S_0 = 720.65$. We use
nine expirations, with maturities $\tau$ from $0.10$ to $0.88$ years. For each expiration the
risk free rate, the dividend yield and the implied volatility come from the pipeline of
Appendix~\ref{app:impl_vol}: the Treasury curve for $r$, put call parity for $q$, and the out
of the money composite for the smile. Each smile is restricted to the moneyness band
$0.8 \le K/S_0 \le 1.2$, where the quotes are liquid. We keep only expirations with at least
twenty five strikes in this band: a sparser expiration is under-determined, and when it sits a
few days from a well-populated one it contributes a redundant thin slice that the fit cannot
resolve and that simply relaxes to a noisy prior.

The surface is held as one spatial slice per maturity, which is the
parametrization~\eqref{eq:coefficient_slices} of Section~\ref{sec:time_dependent_a}: we do not
fit the maturities one at a time nor interpolate between their smiles. The forward equation is
marched once from $0$ to the last maturity, the coefficient taking the value $a_k$ on the $k$th
interval, and the whole time trajectory of the solution is kept and read off at each maturity.
Every observed price is therefore a sample of one and the same evolving solution, and the
unknown is the stack of slices $(a_1,\dots,a_K)$ solved for jointly, at the cost of one forward
solve and one backward adjoint sweep per iteration whatever the number of slices. The
functional minimized is the space and time version~\eqref{eq:surface_functional} announced
there, with $\Omega_j$ the band of strikes quoted at maturity $\tau_j$, the $H^1$ term
smoothing each slice in strike and the $L_2$ term penalizing the jump between consecutive
slices. The triangular
structure noted in Section~\ref{sec:time_dependent_a} is what makes this well posed maturity by
maturity: the earliest expiration constrains the first slice, the next one the second, and so
on. The prior $a^*_k$ of each slice deserves care.
The instinct is to seed it with that maturity's implied variance $\tfrac12\sigma_{\mathrm{iv}}^2$,
but the $H^1$ term penalizes the difference of \emph{slopes} $\partial_y(a_k - a^*_k)$, and by
the factor of two derived below (Section~\ref{sec:factor_two}) the implied skew is only half the
local skew. Such a prior would anchor the reconstruction at the wrong slope and flatten the
recovered skew toward the implied one wherever the data are weak. We instead seed $a^*_k$ with
the local volatility recovered from the smile by inverting the short maturity relation~\eqref{eq:bbf_harmonic},
which on differentiating $y/\sigma_{\mathrm{iv}} = \int_0^y dx/\sigma_{\mathrm{loc}}$ gives
\begin{equation}\label{eq:bbf_inversion}
    \sigma_{\mathrm{loc}}(y) \;=\; \frac{\sigma_{\mathrm{iv}}(y)}
        {1 - y\,\partial_y\sigma_{\mathrm{iv}}(y)/\sigma_{\mathrm{iv}}(y)}.
\end{equation}
This agrees with the smile at the money and doubles its skew there, consistent
with~\eqref{eq:factor_two}; outside the band it is flattened. Interest and dividend enter as
piecewise constant forward rates recovered from the zero rates.

Figure~\ref{fig:egger_surface} shows the result for $\beta_y = 0.1$ and $\beta_\tau = 1$. The
recovered surface has the shape one expects: a pronounced negative skew, with local volatility
close to $0.15$ at the money, near $0.12$ on the upside, and rising sharply on the downside, to
about $0.40$ at the long maturities and much steeper at the short end where the skew is
strongest. Repricing the quotes at the calibrated surface gives a median relative error of
$0.14\%$ across the nine maturities, with the maximum over liquid strikes staying near $1\%$ at
every expiration. This is markedly tighter than the naive prior $a^* = \tfrac12\sigma_{\mathrm{iv}}^2$,
which gives a median of $0.32\%$ and maxima near $2\%$: that prior anchors the skew at half its
correct value, so wherever the quotes carry little information the $H^1$ penalty pulls the
reconstruction off the prices, and the theory-consistent prior~\eqref{eq:bbf_inversion} removes
the bias. As in the synthetic examples, the surface is well determined near the money and
on the downside, where option prices are large and informative, and relaxes towards the prior
on the upside wing, where out of the money call prices are close to zero and constrain the
coefficient only weakly. This is the same ill-posedness of Section~\ref{sec:convergence_rates},
now seen in real data.

The calibrated surface is free of arbitrage, and here the direct approach has a structural
advantage. Because it reproduces the quotes by solving the forward Dupire equation with a local
variance the optimizer keeps strictly positive ($a \ge 10^{-8}$), the output is the price
surface of a genuine diffusion, so the strike and calendar arbitrages of
Section~\ref{sec:no_arbitrage} cannot even be represented. This is unlike the implied volatility
route of Appendix~\ref{app:impl_vol}, where a fitted surface with the slightest butterfly or
calendar violation feeds a negative local variance into the Dupire formula and the
no-arbitrage conditions must be imposed by hand. We verify that the discrete surface inherits
this: the recovered local volatility stays above $0.08$ everywhere, the risk-neutral density
$\partial^2_{KK}\mathcal{C}$ is non-negative across the observed band at every maturity, and the
total implied variance $\sigma_{\mathrm{iv}}^2\tau$ is strictly increasing in maturity at every
moneyness (Figure~\ref{fig:egger_arbitrage}), the two conditions Bergomi records for a local
volatility surface~\cite{bergomi_stochastic_vol}.

\begin{figure}[h]
    \centering
    \includegraphics[width=\textwidth]{egger_example7_arbitrage}
    \caption{No-arbitrage check of the calibrated surface of Figure~\ref{fig:egger_surface}.
             Left: the risk-neutral density $\partial^2 \mathcal{C}/\partial K^2$ per maturity,
             non-negative across the observed band (no strike, or butterfly, arbitrage). Right:
             the total implied variance $\sigma_{\mathrm{iv}}^2\tau$ against maturity at fixed
             log-moneyness, increasing everywhere (no time, or calendar, arbitrage).}
    \label{fig:egger_arbitrage}
\end{figure}

One feature of the fit points to the next section. Near the money the recovered local
volatility is visibly steeper in strike than the market implied volatility plotted on the same
axes. This is a definite relation between the two skews, which we make precise next following
Derman and Bergomi.

\begin{figure}[h]
    \centering
    \includegraphics[width=\textwidth]{egger_example4_surface}
    \caption{Local volatility surface calibrated from SPY quotes on 30 April 2026
             ($S_0 = 720.65$, nine maturities), with $\beta_y = 0.1$ and $\beta_\tau = 1$.
             Left: the recovered surface $\sigma(y,\tau)$ as a 3D mesh. Middle: the recovered local
             volatility slices (solid) against each maturity's market implied volatility
             (faint); the local volatility skew is steeper than the implied one. Right: per
             maturity relative reprice error over liquid strikes (market price at least $0.5\%$
             of $S_0$), median and maximum.}
    \label{fig:egger_surface}
\end{figure}

\section{The local volatility skew and the factor of two}\label{sec:factor_two}

The steeper skew seen in Figure~\ref{fig:egger_surface} is not an artefact of the
calibration. It reflects a general relation between the two objects. The implied volatility
$\sigma_{\mathrm{iv}}(K,\tau)$ is the single Black-Scholes number that reprices the quote at
$(K,\tau)$, while the local volatility $\sigma_{\mathrm{loc}}(S,t)$ is the instantaneous
diffusion coefficient the underlying feels at $(S,t)$. The former is an average of the latter
over the paths that reach the strike, so the smile is a smoothed version of the local
volatility and its skew is the flatter of the two. Derman makes this precise near the
money~\cite{derman_vol_smile}.

The clean statement is the short maturity asymptotic of Berestycki, Busca and
Florent~\cite{berestycki_busca_florent}: as $\tau \to 0$ the implied volatility is the harmonic
average of the local volatility over the log-moneyness interval between spot and strike,
\begin{equation}\label{eq:bbf_harmonic}
    \frac{1}{\sigma_{\mathrm{iv}}(y,\tau)}
      \;=\; \frac{1}{y}\int_0^y \frac{dx}{\sigma_{\mathrm{loc}}(x)} \;+\; o(1),
    \qquad y = \log(K/S_0),
\end{equation}
where $\sigma_{\mathrm{loc}}(x)$ is the local volatility at log-moneyness $x$. Writing
$\sigma_{\mathrm{loc}}(x) = \sigma_0 + \sigma_0' x + O(x^2)$ with $\sigma_0 =
\sigma_{\mathrm{loc}}(0)$ and $\sigma_0' = \partial_y \sigma_{\mathrm{loc}}(0)$, the integrand
is $\sigma_0^{-1}(1 - (\sigma_0'/\sigma_0)\,x) + O(x^2)$, so its average over $[0,y]$ equals
$\sigma_0^{-1}(1 - (\sigma_0'/2\sigma_0)\,y) + O(y^2)$. Inverting~\eqref{eq:bbf_harmonic} gives
$\sigma_{\mathrm{iv}}(y) = \sigma_0 + \tfrac12 \sigma_0'\, y + O(y^2)$, that is
\begin{equation}\label{eq:factor_two}
    \left.\partial_y \sigma_{\mathrm{loc}}\right|_{y=0}
      \;=\; 2\,\left.\partial_y \sigma_{\mathrm{iv}}\right|_{y=0}.
\end{equation}
This is Derman's factor of two: at the money the local volatility skew is twice the implied
volatility skew. Bergomi records the same factor on the dynamic side~\cite{bergomi_stochastic_vol}:
in a local volatility model the at-the-money implied volatility moves with the spot at twice the
rate at which it moves with the strike, so the skew stickiness ratio equals two in the short
maturity limit. The static and the dynamic statements are two faces of the same averaging.

We check~\eqref{eq:factor_two} directly on the calibrated surface. For each maturity we fit a
quadratic to the recovered local volatility slice and to the market implied volatility smile in
a narrow window $|y| \le 0.05$ around the money and read the slope at $y = 0$.
Figure~\ref{fig:egger_skew} shows the two skews and their ratio. Near the money the local skew
is about twice the implied skew, with a mean ratio of $2.15$ for maturities below three months.

One objection must be met. The prior~\eqref{eq:bbf_inversion} already builds the factor of two
into the reconstruction, so the agreement could be an echo of the prior rather than a property
of the data. It is not. Rerunning the calibration with the naive prior
$a^* = \tfrac12\sigma_{\mathrm{iv}}^2$, which instead biases the ratio towards one, leaves the
short maturity ratio at $2.17$, essentially unchanged. The factor of two near the money is
recovered from the option prices themselves and is robust to the choice of prior, which is the
real check that the local volatility, and not just the fitted prices, is meaningful. At long
maturities the ratio is prior sensitive, since the quotes thin out and the reconstruction leans
on the prior there, so we do not read a reliable term structure of the ratio from it; this is
consistent in any case with~\eqref{eq:bbf_harmonic} being a short maturity statement.

\begin{figure}[h]
    \centering
    \includegraphics[width=\textwidth]{egger_example5_skew}
    \caption{The Derman factor of two on the SPY surface of Figure~\ref{fig:egger_surface}.
             Left: the at-the-money implied volatility skew, the recovered local volatility
             skew, and twice the implied skew. Right: the ratio of local to implied skew against
             the reference value two. Both skews are read from a quadratic fit in the window
             $|y| \le 0.05$.}
    \label{fig:egger_skew}
\end{figure}

\section{Out of sample test: pricing between quoted maturities}\label{sec:holdout}

The two constructions, the calibrated local volatility surface and the interpolated implied
volatility surface, agree on the quoted maturities by design. They differ on a contract that
falls between them, and that difference is the test that tells them apart. We remove one
interior expiration from the data, the December 2026 quotes at $\tau = 0.63$, and price it out
of sample two ways. The local volatility surface is recalibrated on the remaining eight
maturities and the Dupire equation is propagated across the gap, so the price at $\tau = 0.63$
is a genuine sample of the forward solution rather than an interpolation of the smile. The
implied volatility construction interpolates the total variance $\sigma_{\mathrm{iv}}^2\tau$
linearly in maturity at fixed log-moneyness between the bracketing expirations $\tau = 0.38$
and $\tau = 0.88$, and reads the price off Black-Scholes.

Figure~\ref{fig:egger_holdout} shows the outcome. Both reproduce the held-out smile closely.
Over the liquid strikes the median relative price error is $0.9\%$ for the implied volatility
interpolation and $1.4\%$ for the local volatility surface, so for pricing vanilla options
between quoted dates the interpolation is marginally the more accurate. This is expected, since
interpolating the variance is essentially built to reproduce the smile. The two part company in
the wings. On the upside, where call prices fall towards zero and the relative error is most
demanding, the implied volatility interpolation reaches almost $10\%$, while the local
volatility price stays under $5\%$, since it is tied to the forward equation rather than to a
straight line in variance.

The conclusion is a measured one. For interpolating vanilla prices the two methods are
comparable, with a slight edge to the lighter implied volatility construction near the money
and to the local volatility surface in the wings. The case for the local volatility surface is
not that it prices vanillas more accurately, but that it delivers a full arbitrage-free spot
dynamics, an instantaneous volatility at every $(S,t)$, which a static implied volatility
interpolation does not. That dynamical content, the same object behind Bergomi's skew
stickiness of Section~\ref{sec:factor_two}, is what is needed to price path-dependent and
forward-starting contracts, for which no single implied volatility exists to interpolate.

\begin{figure}[h]
    \centering
    \includegraphics[width=\textwidth]{egger_example6_holdout}
    \caption{Out-of-sample pricing of a held-out expiration (December 2026, $\tau = 0.63$),
             removed from the calibration and bracketed by $\tau = 0.38$ and $\tau = 0.88$.
             Left: the held-out market implied volatilities against the local volatility surface
             price and the implied volatility interpolation, both converted back to implied
             volatility. Right: relative price error over liquid strikes. The two agree near the
             money; in the upside wing the local volatility price is the more stable.}
    \label{fig:egger_holdout}
\end{figure}